\section{Constructive rectangular bounds and their limits}\label{sec:spatial-new}
Retain $m\geq1$, $1\leq a\leq2m$, exact mean $m$, \eqref{eq:moment},
and $p=D/2$. All assertions in this section remain analytic.

\subsection{A general threshold construction}
For an interval of length $l$, the identity
$\operatorname{Var}(u)=l^{-1}\sum_{x<y}|u(x)-u(y)|^2$ and path
Cauchy--Schwarz give coefficient at most $l(l-1)/4$ per edge: each
crossing pair has distance at most $l-1$, and at most $l^2/4$ pairs cross
any edge. Product variance tensorization therefore yields, for a rectangle
of maximal side $L$,
\begin{equation}\label{eq:boxvar}
 \sum_R|u-\operatorname{mean}_R u|^2
 \leq \frac{L(L-1)}4\sum_{\text{internal forward edges of }R}|du|^2.
\end{equation}
There is no additional factor $D$ in this tensorization.

Fix $\tau\geq1$. If $\tau<c/4$ use singletons. Otherwise set
$\ell=\min\{n,\lfloor2\sqrt{\tau/c}\rfloor\}$ and start with a grid
whose sides have length at most $\ell$. Isolate each vertex with $a<\tau$
as a singleton. Its current rectangle minus that vertex splits into at most
$2D$ rectangles by the first coordinate differing from that vertex. Delete
empty pieces and preserve earlier singletons. Remaining rectangles have
$a\geq\tau$ and side at most $\ell$. Equation~\eqref{eq:boxvar},
nonnegative omitted edges and normalized-mean orthogonality give
\begin{equation}\label{eq:globalmeans}
 c\|U\|^2\leq\E_a(U)+c\|C_{\mathcal P}U\|^2,
 \quad (C_{\mathcal P}U)_{R,j}=\sqrt{|R|}\operatorname{mean}_R U_j
\end{equation}
for ALL cochains. Singleton residuals vanish, including $n=1$; parallel
edge occurrences for $n=2$ are counted separately, and wrap edges can be
omitted positively. With $B_c=(1+\sqrt c)^D$,
\[
 |\mathcal P|\leq B_c\max\{1,n^D/\tau^p\}+2D\#\{a<\tau\},\quad
 \#\{a<\tau\}\leq M n^D\tau^p/m^p.
\]
The singleton branch satisfies this count since $B_c\tau^{-p}>1$ there.
Taking $\tau=\sqrt m$ yields the weaker general bound
\[
 \dim F\leq h(B_c+2DM)\max\{1,n^D/m^{D/4}\}.
\]

\subsection{Every global threshold can lose the target order}
Let $\gamma=\max\{1,c/4\}$, $n=t^2$, $t\geq2$, $m=\gamma n^2$,
$b=\gamma n$, $S=\{0,\ldots,t-1\}^D$, and $f=n^{-p}$. Set $a=b$ on
$S$, $a=A=(m-fb)/(1-f)$ outside. Then $1\leq a\leq2m$, the mean is
$m$, and the critical moment is less than $2-f$. The target maximum is one.
For any global $\tau\geq1$, consider the preceding algorithm with any
translated initial grid or excision order, retaining grid boundaries and
every excised singleton. If $\tau\leq b$, no vertex is weak and the grid
volume bound gives
\[
 |\mathcal P|\geq n^D/(2\sqrt{b/c})^D
     =(c/(4\gamma))^p n^{D/2},
\]
or the singleton branch is larger. For large $n$ the grid side is less
than $n$. If $\tau>b$, all $t^D=n^{D/2}$ weak vertices remain distinct
singletons. Every such threshold therefore loses the target order.

For the SAME profile retain the whole weak cube and split its complement
by the first coordinate outside it, giving $D+1$ rectangles. On $S$,
$ct^2/(4b)=c/(4\gamma)\leq1$; outside, $A\geq m$ and
$cn^2/(4A)\leq1$. Equation~\eqref{eq:boxvar} proves
\eqref{eq:globalmeans} for this constant-size partition. This is a
mechanism obstruction, not a counterexample to every rectangular partition.

\subsection{Supplied scaled cubes attain the target order}
Suppose additionally that disjoint integer cubes $S_k$ of side $L_k$,
nonwrapping in one common root, satisfy for fixed $\beta\geq1$
\begin{equation}\label{eq:scaled-assumptions}
 \begin{split}
 cL_k(L_k-1)/4\leq a(x)\leq\beta L_k^2 &\quad(x\in S_k),\\
 a(x)\geq cn(n-1)/4 &\quad(x\notin\cup_k S_k).
 \end{split}
\end{equation}
These are extra coefficient and geometric assumptions. Each cluster charges
at least $\beta^{-p}$ to $\sum_xa(x)^{-p}$; thus
$J\leq\beta^p M n^D/m^p$.

The solid-fat-object BSP construction of de Berg~\cite{deberg} has linear
size in fixed dimension. Its first stage uses axis-parallel cuts; its second
uses object facet planes. For solid axis-aligned cubes both stages are
axis-parallel. Replace $S_k=\prod_i\{b_{ki},\ldots,b_{ki}+L_k-1\}$ by
$Q_k=\prod_i[b_{ki}-1/4,b_{ki}+L_k-3/4]$. These are positive-side solid
cubes of fixed fatness, strictly disjoint even for touching integer-cell
interpretations; their integer points are precisely $S_k$ and are interior,
including when $L_k=1$.

Clip BSP leaves to $[-1/2,n-1/2]^D$, assign each cut boundary consistently
to one child, take integer points and delete empty cells. If an assigned
integer point belongs to $S_k$, a ball about it lies in $Q_k$, so the open
leaf also meets $Q_k$. Hence one leaf's integer points belong to at most one
cluster. Refine it with that cube's $2D$ facets, at most $3^D$ subrectangles,
each inside one cluster or outside all clusters. Fragmentation is permitted;
whole-cluster preservation is unnecessary. Consequently, for dimension-only
$A_D\geq1$,
\[
 |\mathcal P|\leq A_D\max\{1,J\},\qquad
 \dim F\leq h A_D(1+2^p\beta^pM)\max\{1,n^D/(1+m)^p\}.
\]
Inside a cluster fragment maximal side is at most $L_k$, outside at most
$n$. Equations~\eqref{eq:scaled-assumptions} and~\eqref{eq:boxvar} prove
\eqref{eq:globalmeans}. Projection separation and a high-budget restriction
are unnecessary. This uses classical geometry; it does not derive the
supplied cubes from the original moment.

\subsection{The exact cube decomposition is not universal}
For fixed $c,\beta$ take arbitrarily large integer $n\geq3$, $m=n$,
$W=\{e_1,-e_1\}$, $f=2/n^D$, with $a=1$ on $W$ and
$a=A=(n-f)/(1-f)$ elsewhere. Then $n<A\leq2n$, mean $a=n$, and the
critical moment is less than $2n^{-p}+1-f\leq2$. Choose also
$n>\max\{\beta,8\beta/c,1+8/c\}$ and $n\geq c$.
Every vertex has $a<cn(n-1)/4$, so any family satisfying
\eqref{eq:scaled-assumptions} must cover zero. A side-one cube containing
zero would require $A\leq\beta<n$. A cube of side $L\geq2$ contains at
least one of $e_1,-e_1$, in every translated root. Its lower bound forces
$L^2\leq8/c$, while its upper bound at zero requires
$A\leq\beta L^2\leq8\beta/c<n<A$. Even allowing overlap cannot repair
this exact decomposition.

The SAME profile nevertheless admits a strong grid with
$\ell=\min\{n,\lfloor2\sqrt{A/c}\rfloor\}$ and two singleton excisions.
Equation~\eqref{eq:globalmeans} holds and
\[
 \dim F\leq h\bigl[2^p(1+\sqrt c)^D+4D\bigr]
                  \max\{1,n^D/(1+m)^p\}.
\]
There are at most $\lceil n/\ell\rceil^D+4D$ blocks. If $\ell<n$,
$\ell\geq\sqrt{A/c}$; if $\ell=n$, there is one grid box. Remaining
rectangles have coefficient $A$ and weak residuals are singletons. The
universal supplied-cube route is refuted. The original unrestricted
rectangular question is neither proved nor refuted.
